\chapter{The Room in Arithmetic}
% OUTLINE Ch3 "The Room in Arithmetic (Godel 1931)". THE FOLK-READING
% FENCE CHAPTER (Beastmaster flagged the whole session: Godel = the most
% abused theorem; the no-name-bridge at its most-needed). Laws: both
% inherited (labeled intuitions adds-no-claim); non-funge. CITATION:
% Godel 1931 (arithmetization / Godel numbering, the diagonal lemma,
% first + second incompleteness); the truth-doesn't-matter clause PAID
% here as INDEPENDENCE. The folk-readings (minds/physics/soul) named and
% refused BY NAME. Zero device appearances. Gate: Kodo (the fence
% hardest).

\begin{intuition}
Hand a bureaucracy a rule and it will find the form the rule forgot. Now
imagine a bureaucracy of arithmetic --- a fixed rulebook that proves
statements about whole numbers, each proof a finite paper trail a clerk
could check --- and ask it about a single sentence, carefully built, that
amounts to saying \emph{this office will never stamp me PROVED}. If the
office ever stamps it proved, the office has proved a falsehood and is not
to be trusted. So if the office is trustworthy, it never stamps the
sentence --- which is exactly what the sentence said. The sentence is
true and the office cannot prove it. The feeling to sit with, before the
machinery: nothing mystical happened. A rulebook rich enough to describe
its own stamping was handed a sentence about its own stamping, and the
diagonal did the rest.
\end{intuition}

This is the first of the four great rooms, and its architect is Kurt
G\"odel, whose 1931 paper is the most consequential single document in
the subject and, not coincidentally, the most misread. The construction
reuses the last two chapters' engine exactly, so the reader already owns
its parts; what G\"odel added was the bridge that let arithmetic play all
the roles at once.

The bridge is arithmetization. Fix a scheme --- G\"odel's own, now called
G\"odel numbering --- that assigns to every symbol, every formula, and
every finite proof a distinct whole number, mechanically, so that
statements \emph{about} formulas and proofs become statements about
numbers, and the arithmetic that was the subject matter becomes also the
language. Once formulas are numbers, the rulebook's own notion of
provability --- ``there exists a proof of the formula with number
$n$'' --- is itself an arithmetical statement about $n$, expressible in
the rulebook's own tongue. The system can now talk about itself without
leaving itself; the list of Chapter 2 is the list of provable formulas,
and the digits along its diagonal are facts about which numbers encode
which proofs.

With that in hand the diagonal lemma --- the fixed-point construction
that is Kleene's recursion theorem wearing arithmetic's clothes --- builds
a sentence $G$ that is provably equivalent, inside the system, to the
statement \emph{$G$ has no proof}. Self-application (the sentence is about
its own provability) and negation (it asserts the \emph{absence} of a
proof): the engine of Chapter 2, run in the house of arithmetic. And the
conclusion is G\"odel's first incompleteness theorem: if the system is
consistent, $G$ is neither provable nor refutable within it --- a true
sentence the system cannot reach. His second theorem drives the point
home by turning the sentence on the system's own trustworthiness: no such
system can prove its own consistency, unless it has none. The room is
built, and it does not empty: strengthen the rulebook to prove $G$ and the
construction runs again on the stronger book, producing a fresh $G$ it
cannot reach. The residue re-forms at every strengthening. It always
shows up.

Now the clause this chapter was assigned to pay, because it is the one the
preface promised here and the one the world most often gets backwards.
\emph{The room's truth does not matter} --- and here is what that
precisely means. The trade's word is \emph{independence}: the sentence
$G$ is independent of the system, meaning the system settles neither $G$
nor its negation, and --- this is the part that unsettles first-time
readers --- one may consistently \emph{add} $G$ as a new axiom, or add its
negation, and either way the resulting system is as consistent as the one
before. There are honest arithmetics in which $G$ holds and honest
arithmetics in which it fails. The sentence's truth is not a fact the
house contains; it is a fork the house leaves open, and the house stands
down either branch. That is why the room's truth-value ``does not
matter'' in the book's sense: the room is a structural feature of the
system --- a place the system cannot decide --- and whether one paints
that place true or false is a choice made outside, not a fact read off
inside. The mechanism built the room; the truth is decor.

\begin{intuition}
The picture for independence: a map with a single unlabeled crossroads.
The map is complete and correct everywhere it commits; at the one
crossroads it simply does not say which way is north, and you may label
it either way and get a consistent, usable map. The unlabeled crossroads
is not a flaw in the cartography --- it is a place the cartography, being
honest, declines to decide. G\"odel's sentence is that crossroads, proved
to exist in every map rich enough to chart its own drawing.
\end{intuition}

And now the fence this chapter exists to hold, because G\"odel's theorem
has been conscripted into more arguments it does not support than any
result in the history of mathematics, and this series' whole discipline
converges on refusing exactly that. Incompleteness is a theorem about
\emph{formal systems of sufficient strength to arithmetize their own
syntax}. It says: such a system cannot be both complete and consistent,
and cannot certify its own consistency. It does not say that human minds
transcend machines; the sentence $G$ is unprovable in the system and
equally unprovable by any consistent finite reasoner using the same
axioms, human or otherwise --- the theorem cuts against every consistent
formal reasoner alike, and grants no faculty an exemption. It does not say
the universe is incomputable, or that physics has a G\"odelian limit;
those are claims about the world, and the theorem is a claim about
provability in a formal calculus, connected to the world by no bridge the
theorem itself supplies. It does not say truth is subjective, that
everything is uncertain, or that some knowledge is forever beyond reach in
the inspirational register these words usually travel in. Each of those is
the diagonal's room read as a message about minds, worlds, or souls --- a
no-name-bridge crossing --- and this book refuses each by the discipline
it has run for four volumes: the room is the name of a structure in a
representation, and a representation is not a mind and is not a world.
What G\"odel proved is exact, immense, and entirely about systems of
signs; the rest is import, and the customs house is closed.

The second great room is the semantic twin of this one. Where G\"odel's
sentence spoke of its own \emph{provability} and found a limit inside the
house, the next construction lets a language reach for its own word
\emph{true} --- and finds that the word cannot live in the house at all.
That is Tarski's room, and the next chapter enters it.
